% !PRINTED 186
\weberpage{186}
%<<CONT
\[
\left|
\begin{array}{l}
F''(x_1, x_1),\ F''(x_1, x_2)\ldots F''(x_1, x_m) \\
F''(x_2, x_1),\ F''(x_2, x_2)\ldots F''(x_2, x_m) \\
\cdots\cdots\cdots\cdots\cdots\cdots \\
F''(x_m, x_1),\ F'''(x_m, x_2)\ldots F''(x_m, x_m) % [sic] printed F''' in the second entry of the last row, inconsistent with F'' elsewhere
\end{array}
\right|
\]
und mit $H'$ die entsprechende Determinante für die transformirte
Function $\Phi$ bezeichnen, nach \S.\,56, (8)

\weq{6}{H' = r^2 H.}

$H$ heisst die \gesperrt{Hesse'sche Determinante}\footnote{Genannt nach \gesperrt{Hesse}, der in geometrischen Untersuchungen, namentlich in der Theorie der Curven dritter Ordnung und in Eliminationsproblemen diese Determinante vielfach anwandte und ihre Bedeutung erkannte.} der Function
$F$. In beiden Formeln (5) und (6) bedeutet $r$ die Substitutionsdeterminante
.

\begin{center}
\S.\,60.\\
\gesperrt{Invarianten und Covarianten.}
\end{center}

Bei der linearen Transformation der homogenen Functionen
treten immer gewisse Functionen der Coëfficienten auf, die, wenn
man die Coëfficienten der gegebenen Form durch die Coëfficienten
der transformirten Form ersetzt, sich nicht anders ändern, als
dass sie einen Factor annehmen, der eine Potenz der Substitutionsdeterminante
 ist. Auch bei gleichzeitiger (simultaner) Transformation
 eines Systems von mehreren Formen kommen solche
Functionen vor, die von den Coëfficienten aller Formen des
Systems abhängen. Solche Functionen sind z. B. bei einem
System von $n$ linearen Formen die Determinante des Systems,
und bei einer quadratischen Form die Determinante dieser Form.

Man nennt diese Functionen \gesperrt{Invarianten der Form}, oder
wenn sie von mehreren Formen abhängen, \gesperrt{simultane Invarianten des Systems}.

Wenn also $I(a)$ eine Function der Coëfficienten $a$ einer
Form $F(x)$ ist und wenn $a'$ die Coëfficienten der transformirten
Form $\Phi(x')$ sind, so heisst $I(a)$ eine Invariante, wenn die identische Relation
%CONT>>
